% Sinh từ cấu trúc notebook; mỗi mục có đủ 6 tiểu mục theo đặc tả (Bài 12). Hình/bảng chèn tự động theo tên fig_b*/tab_b*.

\subsection{12.1\quad Xu thế tất định và xu thế ngẫu nhiên; I(d)}
\textit{Bộ dữ liệu: NI, AP, DP (cả 3)}

\subsubsection*{Nội dung lý thuyết cần kiểm chứng}
\TODO{viết ngắn gọn, chính xác theo bài giảng, có số hiệu Định nghĩa/Định lý/Mệnh đề (xem ô ``Nội dung lý thuyết'' của notebook).}

\subsubsection*{Cách triển khai}
\TODO{áp dụng trên bộ nào, kiểu [D] hay [M] (xem ô ``Cách triển khai'').}

\subsubsection*{Mã lệnh Python}
\TODO{dán đoạn mã quan trọng dạng văn bản trong khung \texttt{lstlisting} (không chụp ảnh).}

\subsubsection*{Kết quả tính toán}
\TODO{chọn những bảng/hình dưới đây cần giữ trong báo cáo (có thể bỏ bớt), sắp xếp và viết câu dẫn.}

\input{\tabdir/tab_b12_01_levels_vs_differences.tex}

\figx{fig_b12_01_levels_vs_differences}{Chuỗi, ACF của chuỗi gốc và ACF của sai phân bậc một cho NI, AP (log, huấn luyện), DP (log)}

\subsubsection*{Phân tích và diễn giải}
\TODO{giải thích ý nghĩa thống kê và liên hệ trực tiếp với lý thuyết.}

\subsubsection*{Nhận xét}
\TODO{phù hợp lý thuyết hay không; nếu lệch thì nguyên nhân hợp lý.}

\subsection{12.2\quad Bước ngẫu nhiên}
\textit{Bộ dữ liệu: DP (mô phỏng, tham số từ DP); NI, AP: giải thích}

\subsubsection*{Nội dung lý thuyết cần kiểm chứng}
\TODO{viết ngắn gọn, chính xác theo bài giảng, có số hiệu Định nghĩa/Định lý/Mệnh đề (xem ô ``Nội dung lý thuyết'' của notebook).}

\subsubsection*{Cách triển khai}
\TODO{áp dụng trên bộ nào, kiểu [D] hay [M] (xem ô ``Cách triển khai'').}

\subsubsection*{Mã lệnh Python}
\TODO{dán đoạn mã quan trọng dạng văn bản trong khung \texttt{lstlisting} (không chụp ảnh).}

\subsubsection*{Kết quả tính toán}
\TODO{chọn những bảng/hình dưới đây cần giữ trong báo cáo (có thể bỏ bớt), sắp xếp và viết câu dẫn.}

\input{\tabdir/tab_b12_02_rw_correlation.tex}

\input{\tabdir/tab_b12_02_rw_sample_acf.tex}

\input{\tabdir/tab_b12_02_rw_variance.tex}

\input{\tabdir/tab_b12_02_spurious_regression.tex}

\figx{fig_b12_02_rw_variance}{Bước ngẫu nhiên mô phỏng: $\operatorname{Var}Y_t$ so với $t\hat\sigma^2$ (trái), 50 đường mẫu (phải)}

\subsubsection*{Phân tích và diễn giải}
\TODO{giải thích ý nghĩa thống kê và liên hệ trực tiếp với lý thuyết.}

\subsubsection*{Nhận xét}
\TODO{phù hợp lý thuyết hay không; nếu lệch thì nguyên nhân hợp lý.}

\subsection{12.3\quad Sai phân sửa được gì và làm hỏng gì: sai phân quá mức}
\textit{Bộ dữ liệu: NI, DP (AP: xem 12.10)}

\subsubsection*{Nội dung lý thuyết cần kiểm chứng}
\TODO{viết ngắn gọn, chính xác theo bài giảng, có số hiệu Định nghĩa/Định lý/Mệnh đề (xem ô ``Nội dung lý thuyết'' của notebook).}

\subsubsection*{Cách triển khai}
\TODO{áp dụng trên bộ nào, kiểu [D] hay [M] (xem ô ``Cách triển khai'').}

\subsubsection*{Mã lệnh Python}
\TODO{dán đoạn mã quan trọng dạng văn bản trong khung \texttt{lstlisting} (không chụp ảnh).}

\subsubsection*{Kết quả tính toán}
\TODO{chọn những bảng/hình dưới đây cần giữ trong báo cáo (có thể bỏ bớt), sắp xếp và viết câu dẫn.}

\input{\tabdir/tab_b12_03_cost_simulation.tex}

\input{\tabdir/tab_b12_03_overdifferencing.tex}

\subsubsection*{Phân tích và diễn giải}
\TODO{giải thích ý nghĩa thống kê và liên hệ trực tiếp với lý thuyết.}

\subsubsection*{Nhận xét}
\TODO{phù hợp lý thuyết hay không; nếu lệch thì nguyên nhân hợp lý.}

\subsection{12.4\quad Kiểm định nghiệm đơn vị: ADF và KPSS}
\textit{Bộ dữ liệu: NI, AP, DP (cả 3) + mô phỏng}

\subsubsection*{Nội dung lý thuyết cần kiểm chứng}
\TODO{viết ngắn gọn, chính xác theo bài giảng, có số hiệu Định nghĩa/Định lý/Mệnh đề (xem ô ``Nội dung lý thuyết'' của notebook).}

\subsubsection*{Cách triển khai}
\TODO{áp dụng trên bộ nào, kiểu [D] hay [M] (xem ô ``Cách triển khai'').}

\subsubsection*{Mã lệnh Python}
\TODO{dán đoạn mã quan trọng dạng văn bản trong khung \texttt{lstlisting} (không chụp ảnh).}

\subsubsection*{Kết quả tính toán}
\TODO{chọn những bảng/hình dưới đây cần giữ trong báo cáo (có thể bỏ bớt), sắp xếp và viết câu dẫn.}

\input{\tabdir/tab_b12_04_adf_power.tex}

\input{\tabdir/tab_b12_04_null_distribution.tex}

\input{\tabdir/tab_b12_04_unit_root_tests.tex}

\figx{fig_b12_04_dickey_fuller_null}{Phân phối mô phỏng của $\tau$ dưới H$_0$ so với $N(0,1)$ và giá trị tới hạn 5\%}

\subsubsection*{Phân tích và diễn giải}
\TODO{giải thích ý nghĩa thống kê và liên hệ trực tiếp với lý thuyết.}

\subsubsection*{Nhận xét}
\TODO{phù hợp lý thuyết hay không; nếu lệch thì nguyên nhân hợp lý.}

\subsection{12.5\quad Mô hình ARIMA(p, d, q) và độ trôi}
\textit{Bộ dữ liệu: NI, DP (AP: xem 12.10)}

\subsubsection*{Nội dung lý thuyết cần kiểm chứng}
\TODO{viết ngắn gọn, chính xác theo bài giảng, có số hiệu Định nghĩa/Định lý/Mệnh đề (xem ô ``Nội dung lý thuyết'' của notebook).}

\subsubsection*{Cách triển khai}
\TODO{áp dụng trên bộ nào, kiểu [D] hay [M] (xem ô ``Cách triển khai'').}

\subsubsection*{Mã lệnh Python}
\TODO{dán đoạn mã quan trọng dạng văn bản trong khung \texttt{lstlisting} (không chụp ảnh).}

\subsubsection*{Kết quả tính toán}
\TODO{chọn những bảng/hình dưới đây cần giữ trong báo cáo (có thể bỏ bớt), sắp xếp và viết câu dẫn.}

\input{\tabdir/tab_b12_05_arima_params.tex}

\input{\tabdir/tab_b12_05_fit_stats.tex}

\input{\tabdir/tab_b12_05_rw_manual_vs_arima.tex}

\figx{fig_b12_05_drift_forecast}{DP: dự báo 100 ngày giao dịch của log DAX có và không có độ trôi}

\subsubsection*{Phân tích và diễn giải}
\TODO{giải thích ý nghĩa thống kê và liên hệ trực tiếp với lý thuyết.}

\subsubsection*{Nhận xét}
\TODO{phù hợp lý thuyết hay không; nếu lệch thì nguyên nhân hợp lý.}

\subsection{12.6\quad Ảnh hưởng vĩnh viễn của cú sốc: trọng số $\Psi$}
\textit{Bộ dữ liệu: NI, DP (AP: giải thích)}

\subsubsection*{Nội dung lý thuyết cần kiểm chứng}
\TODO{viết ngắn gọn, chính xác theo bài giảng, có số hiệu Định nghĩa/Định lý/Mệnh đề (xem ô ``Nội dung lý thuyết'' của notebook).}

\subsubsection*{Cách triển khai}
\TODO{áp dụng trên bộ nào, kiểu [D] hay [M] (xem ô ``Cách triển khai'').}

\subsubsection*{Mã lệnh Python}
\TODO{dán đoạn mã quan trọng dạng văn bản trong khung \texttt{lstlisting} (không chụp ảnh).}

\subsubsection*{Kết quả tính toán}
\TODO{chọn những bảng/hình dưới đây cần giữ trong báo cáo (có thể bỏ bớt), sắp xếp và viết câu dẫn.}

\input{\tabdir/tab_b12_06_psi_weights.tex}

\figx{fig_b12_06_psi_weights}{Ảnh hưởng của một cú sốc đơn vị lên $Y_{t+k}$ ($\Psi_k$)}

\subsubsection*{Phân tích và diễn giải}
\TODO{giải thích ý nghĩa thống kê và liên hệ trực tiếp với lý thuyết.}

\subsubsection*{Nhận xét}
\TODO{phù hợp lý thuyết hay không; nếu lệch thì nguyên nhân hợp lý.}

\subsection{12.7\quad ARIMA(0,1,1) $\equiv$ làm trơn mũ đơn $\equiv$ mô hình mức cục bộ}
\textit{Bộ dữ liệu: NI (chính); DP, AP: giải thích}

\subsubsection*{Nội dung lý thuyết cần kiểm chứng}
\TODO{viết ngắn gọn, chính xác theo bài giảng, có số hiệu Định nghĩa/Định lý/Mệnh đề (xem ô ``Nội dung lý thuyết'' của notebook).}

\subsubsection*{Cách triển khai}
\TODO{áp dụng trên bộ nào, kiểu [D] hay [M] (xem ô ``Cách triển khai'').}

\subsubsection*{Mã lệnh Python}
\TODO{dán đoạn mã quan trọng dạng văn bản trong khung \texttt{lstlisting} (không chụp ảnh).}

\subsubsection*{Kết quả tính toán}
\TODO{chọn những bảng/hình dưới đây cần giữ trong báo cáo (có thể bỏ bớt), sắp xếp và viết câu dẫn.}

\input{\tabdir/tab_b12_07_local_level.tex}

\input{\tabdir/tab_b12_07_ses_forecasts.tex}

\subsubsection*{Phân tích và diễn giải}
\TODO{giải thích ý nghĩa thống kê và liên hệ trực tiếp với lý thuyết.}

\subsubsection*{Nhận xét}
\TODO{phù hợp lý thuyết hay không; nếu lệch thì nguyên nhân hợp lý.}

\subsection{12.8\quad Dự báo h bước tối ưu MSE, phương sai sai số và khoảng dự báo}
\textit{Bộ dữ liệu: NI, AP, DP (cả 3) + mô phỏng độ phủ}

\subsubsection*{Nội dung lý thuyết cần kiểm chứng}
\TODO{viết ngắn gọn, chính xác theo bài giảng, có số hiệu Định nghĩa/Định lý/Mệnh đề (xem ô ``Nội dung lý thuyết'' của notebook).}

\subsubsection*{Cách triển khai}
\TODO{áp dụng trên bộ nào, kiểu [D] hay [M] (xem ô ``Cách triển khai'').}

\subsubsection*{Mã lệnh Python}
\TODO{dán đoạn mã quan trọng dạng văn bản trong khung \texttt{lstlisting} (không chụp ảnh).}

\subsubsection*{Kết quả tính toán}
\TODO{chọn những bảng/hình dưới đây cần giữ trong báo cáo (có thể bỏ bớt), sắp xếp và viết câu dẫn.}

\input{\tabdir/tab_b12_08_coverage.tex}

\input{\tabdir/tab_b12_08_forecast_variance.tex}

\figx{fig_b12_08_ni_forecast}{NI: dự báo ARIMA(0,1,1) và khoảng dự báo 95\% cho 20 năm tiếp theo}

\figx{fig_b12_08_se_growth}{Sai số chuẩn của dự báo $h$ bước ($\sigma=1$): $d=0$ so với $d=1$}

\subsubsection*{Phân tích và diễn giải}
\TODO{giải thích ý nghĩa thống kê và liên hệ trực tiếp với lý thuyết.}

\subsubsection*{Nhận xét}
\TODO{phù hợp lý thuyết hay không; nếu lệch thì nguyên nhân hợp lý.}

\subsection{12.9\quad Dự báo trên thang log: trung vị, trung bình, khoảng chuyển đổi}
\textit{Bộ dữ liệu: AP, DP (NI: không biến đổi)}

\subsubsection*{Nội dung lý thuyết cần kiểm chứng}
\TODO{viết ngắn gọn, chính xác theo bài giảng, có số hiệu Định nghĩa/Định lý/Mệnh đề (xem ô ``Nội dung lý thuyết'' của notebook).}

\subsubsection*{Cách triển khai}
\TODO{áp dụng trên bộ nào, kiểu [D] hay [M] (xem ô ``Cách triển khai'').}

\subsubsection*{Mã lệnh Python}
\TODO{dán đoạn mã quan trọng dạng văn bản trong khung \texttt{lstlisting} (không chụp ảnh).}

\subsubsection*{Kết quả tính toán}
\TODO{chọn những bảng/hình dưới đây cần giữ trong báo cáo (có thể bỏ bớt), sắp xếp và viết câu dẫn.}

\input{\tabdir/tab_b12_09_ap_backtransform.tex}

\input{\tabdir/tab_b12_09_dp_forecast.tex}

\figx{fig_b12_09_ap_forecast}{AP: dự báo airline 1959--1960 trên thang gốc (trung vị, khoảng 95\%) và giá trị thực}

\subsubsection*{Phân tích và diễn giải}
\TODO{giải thích ý nghĩa thống kê và liên hệ trực tiếp với lý thuyết.}

\subsubsection*{Nhận xét}
\TODO{phù hợp lý thuyết hay không; nếu lệch thì nguyên nhân hợp lý.}

\subsection{12.10\quad Sai phân mùa vụ, SARIMA và mô hình airline}
\textit{Bộ dữ liệu: AP}

\subsubsection*{Nội dung lý thuyết cần kiểm chứng}
\TODO{viết ngắn gọn, chính xác theo bài giảng, có số hiệu Định nghĩa/Định lý/Mệnh đề (xem ô ``Nội dung lý thuyết'' của notebook).}

\subsubsection*{Cách triển khai}
\TODO{áp dụng trên bộ nào, kiểu [D] hay [M] (xem ô ``Cách triển khai'').}

\subsubsection*{Mã lệnh Python}
\TODO{dán đoạn mã quan trọng dạng văn bản trong khung \texttt{lstlisting} (không chụp ảnh).}

\subsubsection*{Kết quả tính toán}
\TODO{chọn những bảng/hình dưới đây cần giữ trong báo cáo (có thể bỏ bớt), sắp xếp và viết câu dẫn.}

\input{\tabdir/tab_b12_10_airline_acf.tex}

\input{\tabdir/tab_b12_10_sd.tex}

\input{\tabdir/tab_b12_10_seasonality_models.tex}

\figx{fig_b12_10_airline_acf}{ACF của $\nabla\nabla_{12}$ log AP (que) và ACF lý thuyết của mô hình airline (chấm)}

\subsubsection*{Phân tích và diễn giải}
\TODO{giải thích ý nghĩa thống kê và liên hệ trực tiếp với lý thuyết.}

\subsubsection*{Nhận xét}
\TODO{phù hợp lý thuyết hay không; nếu lệch thì nguyên nhân hợp lý.}

\subsection{12.11\quad Mốc chuẩn và thước đo độ chính xác}
\textit{Bộ dữ liệu: NI, AP, DP (cả 3)}

\subsubsection*{Nội dung lý thuyết cần kiểm chứng}
\TODO{viết ngắn gọn, chính xác theo bài giảng, có số hiệu Định nghĩa/Định lý/Mệnh đề (xem ô ``Nội dung lý thuyết'' của notebook).}

\subsubsection*{Cách triển khai}
\TODO{áp dụng trên bộ nào, kiểu [D] hay [M] (xem ô ``Cách triển khai'').}

\subsubsection*{Mã lệnh Python}
\TODO{dán đoạn mã quan trọng dạng văn bản trong khung \texttt{lstlisting} (không chụp ảnh).}

\subsubsection*{Kết quả tính toán}
\TODO{chọn những bảng/hình dưới đây cần giữ trong báo cáo (có thể bỏ bớt), sắp xếp và viết câu dẫn.}

\input{\tabdir/tab_b12_11_ap_benchmarks.tex}

\input{\tabdir/tab_b12_11_dp_multistep.tex}

\input{\tabdir/tab_b12_11_dp_one_step.tex}

\input{\tabdir/tab_b12_11_ni_benchmarks.tex}

\subsubsection*{Phân tích và diễn giải}
\TODO{giải thích ý nghĩa thống kê và liên hệ trực tiếp với lý thuyết.}

\subsubsection*{Nhận xét}
\TODO{phù hợp lý thuyết hay không; nếu lệch thì nguyên nhân hợp lý.}

\subsection{12.12\quad Thẩm định trên gốc cuốn chiếu và vai trò của bước chọn mô hình}
\textit{Bộ dữ liệu: NI, AP, DP (cả 3)}

\subsubsection*{Nội dung lý thuyết cần kiểm chứng}
\TODO{viết ngắn gọn, chính xác theo bài giảng, có số hiệu Định nghĩa/Định lý/Mệnh đề (xem ô ``Nội dung lý thuyết'' của notebook).}

\subsubsection*{Cách triển khai}
\TODO{áp dụng trên bộ nào, kiểu [D] hay [M] (xem ô ``Cách triển khai'').}

\subsubsection*{Mã lệnh Python}
\TODO{dán đoạn mã quan trọng dạng văn bản trong khung \texttt{lstlisting} (không chụp ảnh).}

\subsubsection*{Kết quả tính toán}
\TODO{chọn những bảng/hình dưới đây cần giữ trong báo cáo (có thể bỏ bớt), sắp xếp và viết câu dẫn.}

\input{\tabdir/tab_b12_12_ap_one_step.tex}

\input{\tabdir/tab_b12_12_ap_rolling_by_horizon.tex}

\input{\tabdir/tab_b12_12_dp_rolling.tex}

\input{\tabdir/tab_b12_12_ni_rolling.tex}

\figx{fig_b12_12_ap_rmse_by_horizon}{AP: RMSE gốc cuốn chiếu theo chân trời, airline so với naive mùa vụ}

\subsubsection*{Phân tích và diễn giải}
\TODO{giải thích ý nghĩa thống kê và liên hệ trực tiếp với lý thuyết.}

\subsubsection*{Nhận xét}
\TODO{phù hợp lý thuyết hay không; nếu lệch thì nguyên nhân hợp lý.}

\subsection{12.13\quad Hàm mất mát xác định phiếm hàm; kiểm định Diebold–Mariano}
\textit{Bộ dữ liệu: AP (chính), NI}

\subsubsection*{Nội dung lý thuyết cần kiểm chứng}
\TODO{viết ngắn gọn, chính xác theo bài giảng, có số hiệu Định nghĩa/Định lý/Mệnh đề (xem ô ``Nội dung lý thuyết'' của notebook).}

\subsubsection*{Cách triển khai}
\TODO{áp dụng trên bộ nào, kiểu [D] hay [M] (xem ô ``Cách triển khai'').}

\subsubsection*{Mã lệnh Python}
\TODO{dán đoạn mã quan trọng dạng văn bản trong khung \texttt{lstlisting} (không chụp ảnh).}

\subsubsection*{Kết quả tính toán}
\TODO{chọn những bảng/hình dưới đây cần giữ trong báo cáo (có thể bỏ bớt), sắp xếp và viết câu dẫn.}

\input{\tabdir/tab_b12_13_diebold_mariano.tex}

\input{\tabdir/tab_b12_13_loss_functional.tex}

\input{\tabdir/tab_b12_13_ranking.tex}

\subsubsection*{Phân tích và diễn giải}
\TODO{giải thích ý nghĩa thống kê và liên hệ trực tiếp với lý thuyết.}

\subsubsection*{Nhận xét}
\TODO{phù hợp lý thuyết hay không; nếu lệch thì nguyên nhân hợp lý.}

\subsection{12.14\quad Checklist báo cáo dự báo ARIMA}
\textit{Bộ dữ liệu: NI, AP, DP (cả 3)}

\subsubsection*{Nội dung lý thuyết cần kiểm chứng}
\TODO{viết ngắn gọn, chính xác theo bài giảng, có số hiệu Định nghĩa/Định lý/Mệnh đề (xem ô ``Nội dung lý thuyết'' của notebook).}

\subsubsection*{Cách triển khai}
\TODO{áp dụng trên bộ nào, kiểu [D] hay [M] (xem ô ``Cách triển khai'').}

\subsubsection*{Mã lệnh Python}
\TODO{dán đoạn mã quan trọng dạng văn bản trong khung \texttt{lstlisting} (không chụp ảnh).}

\subsubsection*{Kết quả tính toán}
\TODO{chọn những bảng/hình dưới đây cần giữ trong báo cáo (có thể bỏ bớt), sắp xếp và viết câu dẫn.}

\input{\tabdir/tab_b12_14_checklist.tex}

\figx{fig_b12_14_residual_acf}{ACF của phần dư (trái) và của phần dư bình phương (phải) cho NI, AP, DP}

\subsubsection*{Phân tích và diễn giải}
\TODO{giải thích ý nghĩa thống kê và liên hệ trực tiếp với lý thuyết.}

\subsubsection*{Nhận xét}
\TODO{phù hợp lý thuyết hay không; nếu lệch thì nguyên nhân hợp lý.}

\subsection{$\star$\quad Mở rộng (tùy chọn): hồi quy giả tạo và đồng tích hợp log DAX $\sim$ log SMI}
\textit{Bộ dữ liệu: DAX và SMI (phần mở rộng, không thuộc 6 bộ chính)}

\subsubsection*{Nội dung lý thuyết cần kiểm chứng}
\TODO{viết ngắn gọn, chính xác theo bài giảng, có số hiệu Định nghĩa/Định lý/Mệnh đề (xem ô ``Nội dung lý thuyết'' của notebook).}

\subsubsection*{Cách triển khai}
\TODO{áp dụng trên bộ nào, kiểu [D] hay [M] (xem ô ``Cách triển khai'').}

\subsubsection*{Mã lệnh Python}
\TODO{dán đoạn mã quan trọng dạng văn bản trong khung \texttt{lstlisting} (không chụp ảnh).}

\subsubsection*{Kết quả tính toán}
\TODO{chọn những bảng/hình dưới đây cần giữ trong báo cáo (có thể bỏ bớt), sắp xếp và viết câu dẫn.}

\input{\tabdir/tab_b12_ext_cointegration.tex}

\subsubsection*{Phân tích và diễn giải}
\TODO{giải thích ý nghĩa thống kê và liên hệ trực tiếp với lý thuyết.}

\subsubsection*{Nhận xét}
\TODO{phù hợp lý thuyết hay không; nếu lệch thì nguyên nhân hợp lý.}
